\newpage

\cleardoublepage
\phantomsection
\addcontentsline{toc}{chapter}{References}
\markboth{References}{References}

\begingroup
\sloppy
\setlength{\emergencystretch}{3em}

\begin{thebibliography}{999}

\bibitem{zhouyi} \textbf{\textit{Zhou Yi} (the \textit{I Ching}, or Book of Changes)}, including the `I Ching' and the `Yi Zhuan'. Standard received text.

The phrase `no plain without a slope, no going without a return', quoted on the cover, is drawn from the line statement of the third line (the ninth-yang line) of the `Tai' hexagram of the \textit{Zhou Yi}. The growth and decline of yin and yang, and the idea of recurrent cycles, are an ancient reverberation of the periodic mathematics programme of `symmetry $\to$ periodicity'.

\bibitem{daodejing} \textbf{\textit{Daodejing} (the \textit{Tao Te Ching})}, by Laozi. Standard received text (the Wang Bi commentary edition).

The phrase `The Dao begot the one; the one begot the two; the two begot the three; and the three begot the myriad things', quoted on the cover, is from chapter 42.

\bibitem{zhuangzi} \textbf{\textit{Zhuangzi}}, by Zhuang Zhou. Standard received text.

`Born together with life, dying together with death; dying together with death, born together with life' (the `Qiwulun', `On the Equality of Things') reveals the mutual transformation of opposites.

\bibitem{huangdineijing} \textbf{\textit{Huangdi Neijing} (the \textit{Yellow Emperor's Inner Classic})}, including the \textit{Suwen} and the \textit{Lingshu}. Standard received text.

The system of the five circuits and six qi is an early model of nested periodic structure.

\bibitem{shangshu} \textbf{\textit{Shangshu} (the \textit{Book of Documents})}, including the `Yaodian'. Standard received text.

`A year comprises three hundred and sixty-six days, and by intercalary months the four seasons are fixed and the year completed'---the earliest official record of a lunisolar calendar.

\bibitem{huainanzi} \textbf{\textit{Huainanzi}}, compiled by Liu An. Standard received text.

The preface cites the `Tianwen xun' (Astronomy) and the `Lanming xun' as the central texts for the establishment of the Xuanyuan calendrical epoch and for the interpretation of N\"uwa mending the heavens as a calendrical act.

\bibitem{shanhaijing} \textbf{\textit{Shanhaijing} (the \textit{Classic of Mountains and Seas})}, anonymous. Standard received text.

The preface cites Jingwei filling the sea and Kuafu chasing the sun, as mythological texts for the eternity of calendrical restoration and for the beginnings of observation.

\bibitem{shiji} \textbf{\textit{Shiji} (the \textit{Records of the Grand Historian})}, by Sima Qian.

The `Tianguan shu' (Treatise on the Heavenly Offices) is a systematic record of ancient astronomical observation. The `Li shu' (Treatise on the Calendar) records that the jia-yin day of the first year of the Taichu era served as the baseline for the derivation of the Xuanyuan calendrical epoch.

\bibitem{shuowen} \textbf{\textit{Shuowen Jiezi}}, by Xu Shen.

The preface cites `Intercalary: the month of the surplus division', as the philological basis for the `one day' left behind by Hou Yi shooting the suns, taken as the intercalary surplus.

\bibitem{zhoubi} \textbf{\textit{Zhoubi Suanjing} (the \textit{Arithmetic Classic of the Gnomon and the Circular Paths of Heaven})}, anonymous.

The `seven concentric rings and six intervals' diagram describes the annual apparent motion of the sun by means of concentric circles.

\bibitem{yueling} \textbf{\textit{Yueling} (the `Monthly Ordinances' of the \textit{Liji}, or \textit{Book of Rites})}, anonymous. Standard received text.

It records month by month the periodic rhythms of celestial phenomena, phenology and farming, and is the earliest systematic literature on phenology.

\bibitem{mengzi} \textbf{\textit{Mengzi} (the \textit{Mencius})}, by Meng Ke. Standard received text.

`All things are already complete within me; to turn inward and realise this is joy beyond compare' (the `Jin Xin Shang', `On the Full Development of the Mind, I').

\bibitem{sunzi} \textbf{\textit{Sunzi Bingfa} (the \textit{Art of War})}, by Sun Wu. Standard received text.

`The skilful warrior seeks victory from the configuration of forces, and does not demand it of his men' (the `Bingshi', `On Military Dispositions').

\bibitem{shennong} \textbf{\textit{Shennong Bencao Jing} (Shennong's \textit{Classic of Materia Medica})}, anonymous. Standard received text (the recension edited by Sun Xingyan and Sun Fengyi).

The earliest classic of Chinese materia medica, recording the harvesting seasons and the places of origin of 365 medicinals.

\bibitem{bencao} \textbf{\textit{Bencao Gangmu} (the \textit{Compendium of Materia Medica})}, by Li Shizhen (1596).

Recording 1,892 medicinals, it systematically documents the `season and month', the `locality' and the `processing' of each drug.

\bibitem{huangjijingshi} \textbf{\textit{Huangji Jingshi}}, by Shao Yong.

Its nested periodic system of yuan, hui, yun and shi is the theoretical basis for the nested periodic structure in the `Periodic Structure of Calendar Systems' chapter of this book.

\bibitem{leibniz1703} \textbf{Gottfried Wilhelm Leibniz}. \textit{Explication de l'Arithm\'etique Binaire} (1703).

Leibniz's original paper on binary arithmetic. The discovery of periodicity in binary columns (Periodus numerorum) provides the historical precursor for the bit-width and bandwidth concepts of this book.

\bibitem{leibniz1679} \textbf{Gottfried Wilhelm Leibniz}. \textit{De Progressione Dyadica} (1679).

Leibniz's manuscript on the dyadic progression. The law of periodic variation of each digit in the sequence of natural numbers was first systematically recorded in this manuscript.

\bibitem{leibniz1714} \textbf{Gottfried Wilhelm Leibniz}. \textit{Monadologie} (1714).

The Monadology. The idea that binary arithmetic generates `being' ($1$) from `nothing' ($0$) shares a common philosophical lineage with the Monadology's notion of `constructing a complex world from simple substances'.

\bibitem{galileo1623} \textbf{Galileo Galilei}. \textit{The Assayer} (1623).

The phrase quoted on the cover, `Mathematics is the language with which God has written the universe', originates from this work.

\bibitem{newton1687} \textbf{Isaac Newton}. \textit{Philosophiae Naturalis Principia Mathematica} (1687).

The law of universal gravitation and the three laws of motion provided the first complete mathematical description of classical periodic motion.

\bibitem{euler1748} \textbf{Leonhard Euler}. \textit{Introductio in Analysin Infinitorum} (1748).

Euler's formula $e^{i\theta} = \cos\theta + i\sin\theta$ is the mathematical cornerstone of the epicycle model and of Fourier analysis.

\bibitem{fourier1822} \textbf{Joseph Fourier}. \textit{Th\'eorie Analytique de la Chaleur} (1822).

The original work on Fourier series and the Fourier transform.

\bibitem{hamilton1853} \textbf{William R. Hamilton}. \textit{Lectures on Quaternions}. Hodges and Smith, 1853.

The original work on quaternions.

\bibitem{riemann1854} \textbf{Bernhard Riemann}. \textit{\"Uber die Hypothesen, welche der Geometrie zu Grunde liegen} (1854).

The foundational lecture on Riemannian geometry.

\bibitem{maxwell1873} \textbf{James Clerk Maxwell}. \textit{A Treatise on Electricity and Magnetism} (1873).

Maxwell's equations are the physical prototype of the electromagnetic field theory in the `Periodic Structure of Electromagnetic Systems' chapter of this book.

\bibitem{klein1872} \textbf{Felix Klein}. \textit{Vergleichende Betrachtungen \"uber neuere geometrische Forschungen} (Erlangen Programme) (1872).

Klein's Erlangen programme---geometry is the science of studying the invariants of transformation groups.

\bibitem{poincare1892} \textbf{Henri Poincar\'e}. \textit{Les M\'ethodes Nouvelles de la M\'ecanique C\'eleste} (1892--1899).

Poincar\'e's pioneering study of periodic orbits in the three-body problem.

\bibitem{planck1901} \textbf{Max Planck}. \textit{\"Uber das Gesetz der Energieverteilung im Normalspectrum}. \textit{Annalen der Physik}, 309(3), 553--563 (1901).

The original paper on Planck's formula for black-body radiation.

\bibitem{noether1918} \textbf{Emmy Noether}. \textit{Invariante Variationsprobleme}. \textit{Nachrichten von der Gesellschaft der Wissenschaften zu G\"ottingen}, 235--257 (1918).

The original paper on Noether's theorem. Every continuous symmetry corresponds to a conserved quantity---the theoretical source for the relationship between conservation laws and symmetry in the `Operations and Periodicity' chapter of this book.

\bibitem{einstein1916} \textbf{Albert Einstein}. \textit{Die Grundlage der allgemeinen Relativit\"atstheorie}. \textit{Annalen der Physik}, 49, 769--822 (1916).

The idea that local symmetry determines the dynamics of gravitation.

\bibitem{compton1923} \textbf{Arthur H. Compton}. \textit{A Quantum Theory of the Scattering of X-rays by Light Elements}. \textit{Physical Review}, 21(5), 483--502 (1923).

The original paper on Compton scattering.

\bibitem{yangmills1954} \textbf{Chen Ning Yang, Robert L. Mills}. \textit{Conservation of Isotopic Spin and Isotopic Gauge Invariance}. \textit{Physical Review}, 96(1), 191--195 (1954).

The pioneering paper on Yang--Mills gauge field theory.

\bibitem{weyl1928} \textbf{Hermann Weyl}. \textit{The Theory of Groups and Quantum Mechanics} (1928).

The first systematic application of group representation theory in quantum mechanics.

\bibitem{weyl1950} \textbf{Hermann Weyl}. \textit{The Theory of Groups and Quantum Mechanics}. Dover Publications, 1950.

The English translation of Weyl's classic work on group representation theory. The book's systematic treatment of the representation theory of compact groups, of character theory and of the symmetric group provides the classical framework of representation theory for the `Groups and Group Functions' chapter and the `Group Objects and Higher-order Mappings' chapter of this book.

\bibitem{serre1977} \textbf{Jean-Pierre Serre}. \textit{Linear Representations of Finite Groups}. Springer, 1977.

A classic textbook on group representation theory. Character theory for finite groups and the theory of induced representations are used systematically in the `Periodic Functions and Groups with Periodic Orbits' chapter of this book.

\bibitem{mackey1968} \textbf{George W. Mackey}. \textit{Induced Representations of Groups and Quantum Mechanics}. Benjamin, 1968.

The theory of induced representations. Mackey's theory of induced representations provides the algebraic basis for the functorial formulation of group actions in the `Periodic Category' chapter of this book.

\bibitem{wigner1959} \textbf{Eugene P. Wigner}. \textit{Group Theory and Its Application to the Quantum Mechanics of Atomic Spectra}. Academic Press, 1959.

The Wigner D-matrix is the central tool for harmonic analysis on SU(2) in this book.

\bibitem{vonneumann1955} \textbf{John von Neumann}. \textit{Mathematical Foundations of Quantum Mechanics}. Princeton University Press, 1955.

The Hilbert-space formulation of quantum mechanics. The book's systematic treatment of the spectral theorem for self-adjoint operators and of unitary groups provides the classical operator-theoretic foundation for the spectral theorem in the `Linear Algebra' chapter and for the `Periodic Category' chapter of this book.

\bibitem{dirac1958} \textbf{Paul A. M. Dirac}. \textit{The Principles of Quantum Mechanics}. 4th ed., Oxford University Press, 1958.

Dirac's transformation theory.

\bibitem{weinstein1978} \textbf{Alan Weinstein}. \textit{Periodic orbits for convex Hamiltonian systems}. Annals of Mathematics, 108(3):507--518, 1978.
\bibitem{arnold1989} \textbf{Vladimir I. Arnold}. \textit{Mathematical Methods of Classical Mechanics}. 2nd ed., Springer, 1989.

The principle of least action, symplectic geometry, integrable systems and KAM theory. Arnold's systematic treatment of the Liouville--Arnold theorem and of invariant tori is a central reference for the `Periodic Manifold' chapter of this book.

\bibitem{arnold1988} \textbf{V. I. Arnold, V. V. Kozlov, A. I. Neishtadt}. \textit{Dynamical Systems III: Mathematical Aspects of Classical and Celestial Mechanics}. Springer, 1988.

A systematic exposition of KAM theory and of the perturbation theory of periodic orbits; an important reference for the `Operations and Periodicity' chapter and the `Periodic Structure of Astronomical Systems' chapter of this book.

\bibitem{meyer1992} \textbf{Kenneth R. Meyer, Glen R. Hall}. \textit{Introduction to Hamiltonian Dynamical Systems and the N-Body Problem}. Springer, 1992.

A classic textbook on Hamiltonian dynamical systems, the continuation of periodic orbits, and KAM theory.

\bibitem{perko1991} \textbf{Lawrence Perko}. \textit{Differential Equations and Dynamical Systems}. Springer, 1991.

Its systematic treatment of periodic orbits, limit cycles, Poincar\'e maps and the Poincar\'e--Bendixson theorem in the global-theory chapter provides a standard reference for the closed-orbit criterion in the `Operations and Periodicity' chapter and for the limit-cycle analysis in the `Periodic Structures in Biological Systems' chapter of this book.

\bibitem{iooss1990} \textbf{G\'erard Iooss, Daniel D. Joseph}. \textit{Elementary Stability and Bifurcation Theory}. 2nd ed., Springer, 1990.

A classic textbook on Hopf bifurcation and the branching of periodic solutions. The book's discussion of the branching from steady states to periodic solutions in the two-dimensional and the general case, of subharmonic branching and of quasi-periodic solutions provides the bifurcation-theoretic foundation for the generation mechanism of periodic solutions in the `Operations and Periodicity' chapter and for the emergence of oscillations in the `Periodic Structures in Biological Systems' chapter of this book.

\bibitem{chicone2006} \textbf{Carmen Chicone}. \textit{Ordinary Differential Equations with Applications}. 2nd ed., Springer, 2006.

Ordinary differential equations and their applications, including a systematic treatment of Floquet theory and periodic solutions.

\bibitem{hirsch2013} \textbf{Morris W. Hirsch, Stephen Smale, Robert L. Devaney}. \textit{Differential Equations, Dynamical Systems, and an Introduction to Chaos}. 3rd ed., Academic Press, 2013.

A classic textbook on differential equations, dynamical systems and chaos, including periodic orbits and limit sets.

\bibitem{ruelle1989} \textbf{David Ruelle}. \textit{Elements of Differentiable Dynamics and Bifurcation Theory}. Academic Press, 1989.

Differentiable dynamical systems, fixed points and periodic orbits, stable and unstable manifolds.

\bibitem{robinson2004} \textbf{R. Clark Robinson}. \textit{An Introduction to Dynamical Systems: Continuous and Discrete}. Pearson Prentice Hall, 2004.

An introductory textbook on continuous and discrete dynamical systems, including periodic orbits and chaotic attractors.

\bibitem{strogatz2015} \textbf{Steven H. Strogatz}. \textit{Nonlinear Dynamics and Chaos: With Applications to Physics, Biology, Chemistry, and Engineering}. 2nd ed., Westview Press, 2015.

A classic textbook on nonlinear dynamics and chaos; the standard reference for limit cycles and periodic oscillations.

\bibitem{kuchment1993} \textbf{Peter Kuchment}. \textit{Floquet Theory for Partial Differential Equations}. Birkh\"auser, 1993.

The authoritative work on Floquet theory for partial differential equations. The book's systematic treatment of the Floquet--Lyapunov theorem for linear differential equations with periodic coefficients, of the spectral theory of periodic differential operators, and of the applications of periodic partial differential equations in quantum mechanics, fluid mechanics and elasticity theory constitutes a central reference for Floquet theory in the `Differential Equations and Analysis' chapter of this book.

\bibitem{brown2013} \textbf{B. Malcolm Brown, Michael S. P. Eastham, Karl Michael Schmidt}. \textit{Periodic Differential Operators}. Springer, 2013.

Floquet theory, oscillation and spectral analysis for periodic differential operators. The book's complete treatment of one-dimensional periodic systems, from Floquet theory to spectral analysis, provides a standard reference for the Bloch-wave theory of differential equations with periodic coefficients and for the `Periodic Manifold' chapter of this book.

\bibitem{steinweiss1971} \textbf{Elias M. Stein, Guido Weiss}. \textit{Introduction to Fourier Analysis on Euclidean Spaces}. Princeton University Press, 1971.

The standard reference for harmonic analysis on Euclidean spaces.

\bibitem{folland1995} \textbf{Gerald B. Folland}. \textit{A Course in Abstract Harmonic Analysis}. CRC Press, 1995.

The theory of harmonic analysis on locally compact groups. Folland's systematic treatment of Pontryagin duality, the Peter--Weyl theorem and the Fourier transform on groups provides a standard reference for the unified framework of the `Harmonic Analysis' chapter of this book.

\bibitem{rudin1962} \textbf{Walter Rudin}. \textit{Fourier Analysis on Groups}. Interscience, 1962.

A classic work on harmonic analysis on locally compact Abelian groups.

\bibitem{katznelson2004} \textbf{Yitzhak Katznelson}. \textit{An Introduction to Harmonic Analysis}. 3rd ed., Cambridge University Press, 2004.

An introductory textbook on harmonic analysis. It includes the Fourier series of periodic functions, $L^2$ convergence and uniform convergence.

\bibitem{zygmund1959} \textbf{Antoni Zygmund}. \textit{Trigonometric Series}. 2nd ed., Cambridge University Press, 1959.

A classic masterpiece on the theory of trigonometric series.

\bibitem{grafakos2009} \textbf{Loukas Grafakos}. \textit{Modern Fourier Analysis}. 2nd ed., Springer, 2009.

A systematic textbook on modern Fourier analysis, including the historical origins of the multivariate trigonometric expansion of periodic functions.

\bibitem{varadhan2000} \textbf{S. R. S. Varadhan}. \textit{Harmonic Analysis}. Courant Institute, 2000.

Introductory lecture notes on harmonic analysis, taking the decomposition of periodic functions into combinations of sines and cosines as its central example.

\bibitem{pontryagin1939} \textbf{Lev S. Pontryagin}. \textit{Topological Groups}. Princeton University Press, 1939.

The original work on the Pontryagin duality theorem. The theoretical foundation for the Pontryagin duality theorem in the `Harmonic Analysis' chapter of this book.

\bibitem{peterweyl1927} \textbf{Fritz Peter, Hermann Weyl}. \textit{Die Vollst\"andigkeit der primitiven Darstellungen einer geschlossenen kontinuierlichen Gruppe}. \textit{Mathematische Annalen}, 97(1), 737--755 (1927).

The original paper on the Peter--Weyl theorem.

\bibitem{fejer1904} \textbf{Leopold Fej\'er}. \textit{Untersuchungen \"uber Fouriersche Reihen}. \textit{Mathematische Annalen}, 58, 51--69 (1904).

The original paper on the Fej\'er kernel.

\bibitem{bohr1925} \textbf{Harald Bohr}. \textit{Zur Theorie der fastperiodischen Funktionen}. \textit{Acta Mathematica}, 45, 29--127 (1925).

The pioneering paper on almost periodic functions; the theoretical source of the concept of almost periodicity in the `Periodic Functions and Groups with Periodic Orbits' chapter of this book.

\bibitem{corduneanu1989} \textbf{C. Corduneanu}. \textit{Almost Periodic Functions}. 2nd English ed., Chelsea, 1989.

A systematic textbook on almost periodic functions, including the classical treatment of the Bohr mean and of the structure of the group with periodic orbits.

\bibitem{terras1999} \textbf{Audrey Terras}. \textit{Fourier Analysis on Finite Groups and Applications}. Cambridge University Press, 1999.

An introductory textbook on Fourier analysis on finite groups; it provides the group-representation-theoretic basis for the phase operators and the Fourier matrices on the cyclic group $C_T$ in this book.

\bibitem{luong2009} \textbf{Bao Luong}. \textit{Fourier Analysis on Finite Abelian Groups}. Birkh\"auser, 2009.

Fourier analysis on finite Abelian groups, including the Fourier transform and the inversion formula for periodic functions.

\bibitem{shannon1948} \textbf{Claude E. Shannon}. \textit{A Mathematical Theory of Communication}. \textit{Bell System Technical Journal}, 27, 379--423, 623--656 (1948).

The foundational work on information entropy and the sampling theorem.

\bibitem{shannon1949} \textbf{Claude E. Shannon}. \textit{Communication in the Presence of Noise}. \textit{Proceedings of the IRE}, 37(1), 10--21 (1949).

The rigorous proof of the Nyquist--Shannon sampling theorem. The theoretical source for the Shannon--Nyquist sampling theorem in the `Periodic Functions and Groups with Periodic Orbits' chapter of this book.

\bibitem{slepian1978} \textbf{David Slepian}. \textit{Prolate Spheroidal Wave Functions, Fourier Analysis, and Uncertainty---V: The Discrete Case}. \textit{Bell System Technical Journal}, 57(5), 1371--1430 (1978).

A classic paper on discrete prolate spheroidal (Slepian) sequences, and the theoretical basis for the discussion of multitaper spectral estimation and energy concentration in this book.

\bibitem{percival1993} \textbf{Donald B. Percival, Andrew T. Walden}. \textit{Spectral Analysis for Physical Applications: Multitaper and Conventional Univariate Techniques}. Cambridge University Press, 1993.

A systematic textbook on multitaper spectral estimation and on classical univariate spectral analysis techniques. The book's complete treatment of Slepian sequences and of Thomson's multitaper method provides a standard reference for the spectral estimation methods in the `Time Series Analysis' chapter and the `Discrete Sampling and Intercalation of Periodic Signals' chapter of this book.

\bibitem{oppenheim2010} \textbf{Alan V. Oppenheim, Ronald W. Schafer}. \textit{Discrete-Time Signal Processing}. 3rd ed., Pearson, 2010.

The standard textbook on discrete-time signal processing, including a detailed treatment of periodic convolution, circular convolution and the DFT.

\bibitem{mallat2009} \textbf{St\'ephane Mallat}. \textit{A Wavelet Tour of Signal Processing}. 3rd ed., Academic Press, 2009.

The authoritative work on wavelet signal processing.

\bibitem{daubechies1992} \textbf{Ingrid Daubechies}. \textit{Ten Lectures on Wavelets}. SIAM, 1992.

A classic textbook on wavelet theory. The Mexican-hat kernel is the standard mother function of the wavelet transform.

\bibitem{cooley1965} \textbf{James W. Cooley, John W. Tukey}. \textit{An Algorithm for the Machine Calculation of Complex Fourier Series}. \textit{Mathematics of Computation}, 19(90), 297--301 (1965).

The original paper on the Cooley--Tukey fast Fourier transform (FFT).

\bibitem{bluestein1970} \textbf{Leo I. Bluestein}. \textit{A Linear Filtering Approach to the Computation of the Discrete Fourier Transform}. \textit{IEEE Transactions on Audio and Electroacoustics}, 18(4), 451--455 (1970).

The original paper on Bluestein's algorithm.

\bibitem{marks1991} \textbf{Robert J. Marks II}. \textit{Introduction to Shannon Sampling and Interpolation Theory}. Springer, 1991.

An introductory textbook on Shannon sampling and interpolation theory, including generalisations of the sampling theorem and cardinal series.

\bibitem{elder2015} \textbf{Yonina C. Eldar}. \textit{Sampling Theory: Beyond Bandlimited Systems}. Cambridge University Press, 2015.

Sampling theory beyond bandlimited systems, including sub-Nyquist sampling and compressed sensing.

\bibitem{khinchin1963} \textbf{A. Ya. Khinchin}. \textit{Continued Fractions}. 3rd ed., P. Noordhoff, 1963.

A classic work on the theory of continued fractions. The book's systematic treatment of the best-approximation property of convergents, of the estimation of the order of approximation and of the metric theory of Diophantine approximation constitutes the mathematical basis for the boundedness theorem for intercalation error and for continued-fraction optimality in the `Discrete Sampling and Intercalation of Periodic Signals' chapter of this book.

\bibitem{lang1995} \textbf{Serge Lang}. \textit{Introduction to Diophantine Approximations}. Springer, 1995.

An introductory textbook on Diophantine approximation, including continued fractions, best approximations and Kronecker's approximation theorem.

\bibitem{walters1981} \textbf{Peter Walters}. \textit{An Introduction to Ergodic Theory}. Springer, 1981.

A classic textbook on ergodic theory. The book's systematic exposition of the Birkhoff ergodic theorem, the Poincar\'e recurrence theorem and the ergodic theory of continuous transformations on compact metric spaces provides a standard reference for ergodicity and for the time-average--space-average criterion in the `Probability and Distribution' chapter of this book.

\bibitem{krengel1985} \textbf{Ulrich Krengel}. \textit{Ergodic Theorems}. de Gruyter, 1985.

The authoritative monograph on ergodic theorems, including systematic proofs of von Neumann's mean ergodic theorem and Birkhoff's pointwise ergodic theorem.

\bibitem{lee2012smooth} \textbf{John M. Lee}. \textit{Introduction to Smooth Manifolds}. 2nd ed., Springer, 2012.

A modern standard textbook on the theory of smooth manifolds. The book's systematic treatment of manifolds, tangent bundles, vector fields and differential forms provides a standard reference for the `Manifold and Fibre Bundle' chapter and the `Periodic Manifold' chapter of this book.

\bibitem{wang2023manifolds} \textbf{\textit{Liuxing yu Xianweicong} (Manifolds and Fibre Bundles)}, by Wang Zhengdong. Publication details to be supplied.

A foundational textbook for the `Manifold and Fibre Bundle' chapter of this book.

\bibitem{nakahara2003} \textbf{Mikio Nakahara}. \textit{Geometry, Topology and Physics}. 2nd ed., IOP Publishing, 2003.

The fibre-bundle formulation of gauge field theory.

\bibitem{conner1979} \textbf{Pierre E. Conner}. \textit{Differentiable Periodic Maps}. Springer, 1979.

The topological theory of differentiable periodic maps; it provides the topological background for periodicity on manifolds in the `Periodic Manifold' chapter of this book.

\bibitem{atiyahsinger1968} \textbf{Michael Atiyah, Isadore M. Singer}. \textit{The Index of Elliptic Operators}. \textit{Annals of Mathematics}, 87(3), 484--530 (1968).

The Atiyah--Singer index theorem.

\bibitem{bohm2003} \textbf{Arno Bohm, Ali Mostafazadeh, Hiroyasu Koizumi, Qian Niu, Joseph Zwanziger}. \textit{The Geometric Phase in Quantum Systems: Foundations, Mathematical Concepts, and Applications in Molecular and Condensed Matter Physics}. Springer, 2003.

A systematic monograph on the geometric phase. The book's systematic discussion of the Berry adiabatic phase, of the geometric connection and curvature, and of the applications of the geometric phase in molecular and condensed-matter physics provides a standard reference for the geometric interpretation of phase in the `Periodic Structure of Electromagnetic Systems' chapter and the `Periodic Manifold' chapter of this book.

\bibitem{feller1950} \textbf{William Feller}. \textit{An Introduction to Probability Theory and Its Applications}, Vol.~1. Wiley, 1950.

A classic textbook on probability theory. The book's systematic treatment of random variables, distribution functions, characteristic functions and limit theorems provides a standard reference for the `Probability and Distribution' chapter of this book.

\bibitem{feller1966} \textbf{William Feller}. \textit{An Introduction to Probability Theory and Its Applications}, Vol.~2. Wiley, 1966.

The second volume of Probability Theory, including a systematic treatment of measure theory, random walks, the Laplace transform and probability distributions.

\bibitem{onsager1944} \textbf{Lars Onsager}. \textit{Crystal Statistics. I. A Two-Dimensional Model with an Order-Disorder Transition}. \textit{Physical Review}, 65(3--4), 117--149 (1944).

Onsager's exact solution of the two-dimensional Ising model, a foundational work in the theory of phase transitions and critical phenomena.

\bibitem{boltzmann1877} \textbf{Ludwig Boltzmann}. \textit{\"Uber die Beziehung zwischen dem zweiten Hauptsatze der mechanischen W\"armetheorie und der Wahrscheinlichkeitsrechnung}. \textit{Wiener Berichte}, 76, 373--435 (1877).

The original paper on Boltzmann's principle $S = k_B \ln W$. The principle links entropy to the geometric area of the constant-energy manifold in phase space, and provides a historical and theoretical foundation for the relationship between entropy and distribution in the `Probability and Distribution' chapter of this book.

\bibitem{landau1987} \textbf{L. D. Landau, E. M. Lifshitz}. \textit{Fluid Mechanics}. 2nd ed., Course of Theoretical Physics, Vol.~6, Pergamon Press, 1987.

A classic textbook on fluid mechanics. The book's systematic treatment of the equations of motion of ideal and viscous fluids, of turbulence and of boundary layers provides a standard reference for the fluid-dynamical background of the `Periodic Structure of Geographic Systems' chapter and the `Periodic Structure of the Climate System' chapter of this book.

\bibitem{arnoldkhesin1998} \textbf{Vladimir I. Arnold, Boris A. Khesin}. \textit{Topological Methods in Hydrodynamics}. Springer, 1998.

The authoritative monograph in the field of topological hydrodynamics. The book treats the topological, group-representation-theoretic and geometric problems of ideal fluids and magnetohydrodynamics from a unified point of view, and provides a standard reference for the topological constraints on fluid motion in the `Periodic Manifold' chapter and the `Periodic Structure of Geographic Systems' chapter of this book.

\bibitem{nocedal2006} \textbf{Jorge Nocedal, Stephen J. Wright}. \textit{Numerical Optimization}. 2nd ed., Springer, 2006.

The standard textbook on numerical optimisation, including a systematic analysis of gradient descent, Newton's method and quasi-Newton methods.

\bibitem{kochenderfer2019} \textbf{Mykel J. Kochenderfer, Tim A. Wheeler}. \textit{Algorithms for Optimization}. MIT Press, 2019.

A practical textbook on optimisation algorithms, including derivatives, gradients and multidimensional optimisation.

\bibitem{boyd2004} \textbf{Stephen P. Boyd, Lieven Vandenberghe}. \textit{Convex Optimization}. Cambridge University Press, 2004.

The standard textbook on convex optimisation.

\bibitem{nesterov2018} \textbf{Yurii Nesterov}. \textit{Lectures on Convex Optimization}. 2nd ed., Springer, 2018.

A classic set of lectures on convex optimisation. Nesterov's accelerated gradient method provides a theoretical reference for the acceleration strategy of multi-window weighted synthesis in KDI.

\bibitem{goodfellow2016} \textbf{Ian Goodfellow, Yoshua Bengio, Aaron Courville}. \textit{Deep Learning}. MIT Press, 2016.

The standard textbook on deep learning, including a systematic discussion of backpropagation and optimisation dynamics.

\bibitem{hirose2006} \textbf{Akira Hirose}. \textit{Complex-Valued Neural Networks}. Springer, 2006.

The first monograph on complex-valued neural networks; it provides the mathematical foundations, in the context of neural networks, for the complex exponential activation function (CEAF) of this book.

\bibitem{hirose2012} \textbf{Akira Hirose} (ed.). \textit{Complex-Valued Neural Networks: Advances and Applications}. Wiley-IEEE Press, 2012.

Advances and applications in complex-valued neural networks, including amplitude--phase activation functions and the complex-valued generalisation of backpropagation.

\bibitem{scholkopf2002} \textbf{Bernhard Sch\"olkopf, Alexander J. Smola}. \textit{Learning with Kernels}. MIT Press, 2002.

The authoritative work on kernel machine learning. The theory of reproducing-kernel Hilbert spaces provides a rigorous foundation for the kernel-function transformation matrix.

\bibitem{jackson1999} \textbf{John D. Jackson}. \textit{Classical Electrodynamics}. 3rd ed., Wiley, 1999.

The authoritative textbook on classical electrodynamics. Its account of the periodic structure of electromagnetic systems, of the derivation of Maxwell's equations and of the Lorentz force formula provides a theoretical reference for this research.

\bibitem{griffiths2013} \textbf{David J. Griffiths}. \textit{Introduction to Electrodynamics}. 4th ed., Pearson, 2013.

The standard introductory textbook on electrodynamics.

\bibitem{alexander2017} \textbf{Charles K. Alexander, Matthew N. O. Sadiku}. \textit{Fundamentals of Electric Circuits}. 6th ed., McGraw-Hill, 2017.

The standard textbook on circuit analysis.

\bibitem{razavi2017} \textbf{Behzad Razavi}. \textit{Design of Analog CMOS Integrated Circuits}. 2nd ed., McGraw-Hill, 2017.

The authoritative textbook on analogue CMOS integrated-circuit design, covering the design theory of crystal oscillators, PLLs and clock trees.

\bibitem{weste2011} \textbf{Neil H. E. Weste, David Money Harris}. \textit{CMOS VLSI Design: A Circuits and Systems Perspective}. 4th ed., Addison-Wesley, 2011.

The standard textbook on VLSI design.

\bibitem{darwin1859} \textbf{Charles Darwin}. \textit{On the Origin of Species}. John Murray, 1859.

The original work on the theory of natural selection.

\bibitem{lotka1925} \textbf{Alfred J. Lotka}. \textit{Elements of Physical Biology}. Williams and Wilkins, 1925.

The original source of the Lotka--Volterra predator--prey equations.

\bibitem{volterra1926} \textbf{Vito Volterra}. \textit{Variazioni e fluttuazioni del numero d'individui in specie animali conviventi}. \textit{Memorie della Regia Accademia Nazionale dei Lincei}, 2, 31--113 (1926).

Volterra derived the predator--prey equations independently of Lotka.

\bibitem{winfree2001} \textbf{Arthur T. Winfree}. \textit{The Geometry of Biological Time}. 2nd ed., Springer, 2001.

The geometric theory of biological time, treating the mathematical structure of rhythmically returning cycles.



A classic textbook on nonlinear dynamics and chaos; the standard reference for limit cycles and periodic oscillations.

\bibitem{aschoff1965} \textbf{J\"urgen Aschoff}. \textit{Circadian Rhythms in Man}. \textit{Science}, 148(3676), 1427--1432 (1965).

A classic study of circadian rhythms in humans.

\bibitem{forger2017} \textbf{Daniel B. Forger}. \textit{Biological Clocks, Rhythms, and Oscillations: The Theory of Biological Timekeeping}. MIT Press, 2017.

A textbook on the mathematical modelling of biological clocks, rhythms and oscillations, including limit-cycle analysis of circadian rhythms and phase response curves.

\bibitem{nowak2006} \textbf{Martin A. Nowak}. \textit{Evolutionary Dynamics}. Harvard University Press, 2006.

The mathematical theory of evolutionary dynamics.

\bibitem{kondratiev1935} \textbf{Nikolai D. Kondratiev}. \textit{The Long Waves in Economic Life}. \textit{Review of Economic Statistics}, 17(6), 105--115 (1935).

Kondratiev's long-wave theory.

\bibitem{kennedy1987} \textbf{Paul Kennedy}. \textit{The Rise and Fall of the Great Powers}. Random House, 1987.

A classic work on the rise and fall of the great powers.

\bibitem{toynbee1934} \textbf{Arnold J. Toynbee}. \textit{A Study of History} (12 vols). Oxford University Press, 1934--1961.

Toynbee's challenge-and-response theory.

\bibitem{braudel1949} \textbf{Fernand Braudel}. \textit{La M\'editerran\'ee et le Monde M\'editerran\'een}. Armand Colin, 1949.

Braudel's multi-timescale view of history.

\bibitem{maclane1998} \textbf{Saunders Mac Lane}. \textit{Categories for the Working Mathematician}. 2nd ed., Springer, 1998.

The standard textbook on category theory; the theoretical basis for functors and natural transformations in the `Periodic Category' chapter of this book.

\bibitem{awodey2010} \textbf{Steve Awodey}. \textit{Category Theory}. 2nd ed., Oxford University Press, 2010.

A modern introductory textbook on category theory.

\bibitem{ren2011} \textbf{Wei Ren, Yongcan Cao}. \textit{Distributed Coordination of Multi-agent Networks}. Springer, 2011.

Distributed coordination of multi-agent networks, including coordination of collective periodic motion and periodic solutions.

\bibitem{mackay2003} \textbf{David J. C. MacKay}. \textit{Information Theory, Inference, and Learning Algorithms}. Cambridge University Press, 2003.

A masterpiece at the intersection of information theory and Bayesian inference.

\bibitem{cover2006} \textbf{Thomas M. Cover, Joy A. Thomas}. \textit{Elements of Information Theory}. 2nd ed., Wiley, 2006.

The standard textbook on information theory.

\bibitem{jeffreys1961} \textbf{Harold Jeffreys}. \textit{Theory of Probability}. 3rd ed., Oxford University Press, 1961.

A classic work on Bayesian statistics.

\bibitem{sutton2018} \textbf{Richard S. Sutton, Andrew G. Barto}. \textit{Reinforcement Learning: An Introduction}. 2nd ed., MIT Press, 2018.

The framework of reward signals and policy optimisation in reinforcement learning.

\bibitem{levinpereswilmer2017} \textbf{David A. Levin, Yuval Peres, Elizabeth L. Wilmer}. \textit{Markov Chains and Mixing Times}. 2nd ed., American Mathematical Society, 2017.

The standard textbook on Markov-chain mixing times; Chapter 12 gives the spectral gap, the contraction of the $\chi^2$ distance and the total-variation bound (Theorem 12.3, Corollary 12.5).

\bibitem{kim2026} \textbf{Eloy Kim}. \textit{Periodicity Mathematics: A Unified Framework of Epicycle Series and Interdisciplinary Periodic Structures}. Manuscript, 2026.

The complete manuscript of this book. It establishes a unified framework for periodic mathematics, comprising the core chapters on periodic foundations, number systems, scalars and vectors, functions and the foundations of analysis, series and convolution, differential equations and complex analysis, linear algebra, manifold foundations and fibre bundles, periodic manifolds, periodic categories, probability theory, vector analysis, groups and group functions, higher-order groups and higher-order functions, periodic functions and groups with periodic orbits, harmonic analysis, time-series analysis and periodic systems.

\end{thebibliography}

\endgroup


\vspace{1cm}
\begin{flushright}
\textit{--- Dedicated to all who seek the `Dao' and the `Number'.}
\end{flushright}

\cleardoublepage
\phantomsection
\addcontentsline{toc}{chapter}{About the Author}
\markboth{About the Author}{About the Author}
\chapter*{About the Author}

Eloy Kim is an independent researcher who has long worked on periodic mathematics, harmonic analysis and interdisciplinary mathematical modelling. The author is proficient in C/C++, Java, TypeScript/JavaScript and Android development, and combines engineering implementation with theoretical derivation. His research interests span the \textit{I Ching}, programming, mathematics and physics, in a sustained attempt to find a unified structure and a common language between philosophy and mathematical science.

\vspace{1em}
\noindent\textbf{Contact:} \texttt{eloy.kim@outlook.com}

\cleardoublepage
\phantomsection
\addcontentsline{toc}{chapter}{Acknowledgements}
\markboth{Acknowledgements}{Acknowledgements}
\chapter*{Acknowledgements}

The successful completion of this research owes a great deal to the valuable research funding and material support provided by many teachers, friends and colleagues. I should like to record my sincere thanks to them. Their support secured the material conditions necessary for the literature review, the mathematical derivations and the interdisciplinary case studies of which this book consists, and enabled the author to concentrate on completing this long-term research programme.

The table below lists the principal funders and supporters (in order of acknowledgement and without any implication of precedence) whose names or institutions are known. The author will continue to add to this table as the programme advances.

\begin{center}
\renewcommand{\arraystretch}{1.35}
\setlength{\tabcolsep}{4pt}
\small
\begin{longtable}{@{}>{\raggedright\arraybackslash}p{1.1cm}
                     >{\raggedright\arraybackslash}p{3.0cm}
                     >{\raggedright\arraybackslash}p{3.4cm}
                     >{\raggedright\arraybackslash}p{5.6cm}@{}}
\caption*{}\\
\toprule
\textbf{No.} & \textbf{Supporter} & \textbf{Type of support} & \textbf{Note of thanks} \\
\midrule
\endfirsthead

\toprule
\textbf{No.} & \textbf{Supporter} & \textbf{Type of support} & \textbf{Note of thanks} \\
\midrule
\endhead

\midrule
\multicolumn{4}{r}{\small continued on the next page} \\
\endfoot

\bottomrule
\endlastfoot

1 & HYC & Material and moral support & Sustained material support and moral encouragement throughout the many years in which this book was written; an essential mainstay without which the research could have been carried on. \\
2 & Tong Su & Research funding & Provided crucial funding and intellectual stimulus at the inception of the research programme, making this research possible. Contact later became less frequent, but that initial impetus was irreplaceable. \\
3 & JD Li & Research funding & Sustained funding of interdisciplinary mathematical exploration, providing stable material support for the literature review and the case studies. \\
4 & An anonymous supporter & Material support & Gave material support in the early stages of this book. The author has been unable to ascertain the name, and records his thanks here. \\
5 & Supporters who left midway & Material support & Once gave financial or material assistance in the early years of the research programme, and subsequently lost contact. The author remembers their contribution and records his gratitude. \\

\end{longtable}
\end{center}

Beyond these, there are a number of friends and institutions who gave financial, documentary, equipment, premises or other material support at various stages of the writing of this book, but whose names the author has not recorded individually, or who prefer not to be named publicly. Their help was equally indispensable, and is sincerely acknowledged here.

Thanks are due to all who have placed their trust in, and given their support to, basic research and interdisciplinary exploration.

\vspace{1em}
\begin{flushright}
The author, Eloy Kim\\
\today
\end{flushright}

\vspace{2cm}
\begin{center}
    \textbf{Periodicity Mathematics}\\
    \copyright\ 2026 \\
    Licence: Apache 2.0
\end{center}

